% =====================================================================
% Response --- Reviewer 2
% Source: docs/reviews/metro-rev2.docx
% Point map: docs/DOCS_ASSESSMENT.md, Stage 3 (R2.1-R2.14).
% Status: STUB --- points restated, responses to be written.
% ---------------------------------------------------------------------
% TONE (Stage 4 decision): maximally gracious. Concede "the bridge exists"
% fully; present the corrected equilibrium and the measured matrix as the
% contribution. Do NOT defend the old bridge-discovery framing.
% =====================================================================
\section*{Reviewer 2}

\noindent General assessment: the IO-CGE bridge already exists (Robinson
2006); the paper's exercise is therefore claimed to be superfluous and to
add no novel insight. Our response concedes the bridge and repositions the
contribution as the measurement of the closure space.

\subsection*{R2.1 --- Labour-supply elasticity as the bridging parameter}
\refcom{[the elasticity of primary factor supply --- the labour supply
elasticity --- is the parameter that governs how far one travels along the
bridge].}
\response TODO: concede; BETA and the matrix; the corrected equilibrium.
\loc Sec. 1; Sec. 4; Sec. 6.

\subsection*{R2.2 --- Closure literature}
\refcom{[the existing literature on alternative factor-market closures in CGE
models].}
\response TODO: the closure taxonomy; Robinson (2006) central.
\loc Sec. 2.

\subsection*{R2.3 --- The 'paradigmatic cleavage' framing}
\refcom{The paper's characterisation of a 'paradigmatic cleavage' \dots\ is
contrived and falls flat \dots\ a 'domain of applicability' angle is more
fruitful.}
\response TODO: reframe via domain of applicability.
\loc Sec. 1.

\subsection*{R2.4 --- The 'miracle' of the closed-economy calibration}
\refcom{[how the closed-economy model was calibrated to Germany's open
economy with its large export and import flows].}
\response TODO: the open-economy pipeline (Sec. 4.1) + the external-debt
closure; the import split; the 5.387 % residual.
\loc Sec. 3; App. B.

\subsection*{R2.5 --- Sector-specific immobile labour}
\refcom{[the alternative, more commonly used in the CGE literature, of labour
as an intersectorally mobile factor with a uniform economy-wide wage].}
\response TODO: the mobile core (ALPHA) is now the default; BF the pole.
\loc Sec. 4; Sec. 6.

\subsection*{R2.6 --- Sign error in eq.\ (19)}
\refcom{Total cost is the sum of --- not the difference between --- the
intermediate input cost term and the primary factor cost term.}
\response TODO: fix.
\loc Sec. 4, eq.\ (cost).

\subsection*{R2.7 --- Supply-vs-demand characterization}
\refcom{In the general equilibrium framework supply and demand are
simultaneously determined --- full stop.}
\response TODO: remove the passage.
\loc Sec. 4 (deleted).

\subsection*{R2.8 --- CES novelty claim}
\refcom{[CES functions have been commonly used in CGE analysis right from
the beginnings].}
\response TODO: drop the novelty claim.
\loc Sec. 1 / Sec. 4.

\subsection*{R2.9 --- Cobb-Douglas as 'the archetypical neoclassical approach'}
\refcom{[the decision that the Cobb-Douglas case is the archetype is entirely
arbitrary].}
\response TODO: drop the archetype claim.
\loc Sec. 4.

\subsection*{R2.10 --- Arbitrary uniform allocation of additional labour}
\refcom{[the entirely arbitrary assumption that 'the additional labor is
transferred uniformly to the sectors'].}
\response TODO: eliminated by endogenous allocation (BF/ALPHA).
\loc Sec. 5; Sec. 6.

\subsection*{R2.11 --- 'Fundamentally contested' allocation claim}
\refcom{[the suggestion that the effects of a demand stimulus on exact
allocation of labor remains fundamentally contested].}
\response TODO: drop or reference.
\loc Sec. 8 (removed).

\subsection*{R2.12 --- Internal mobility contradiction}
\refcom{[section 4.2 argues as if production factors are intersectorally
mobile].}
\response TODO: resolved by the mobile core.
\loc Sec. 4.2.

\subsection*{R2.13 --- McGregor et al.\ (1996)}
\refcom{[see McGregor, Swales \& Yin 1996].}
\response TODO: engage.
\loc Sec. 2; Sec. 6.

\subsection*{R2.14 --- Skill classes}
\refcom{[a proper CGE analysis would split labor by skill class and use
different labor market closures by skill type].}
\response TODO: optional extension; state the limitation.
\loc Sec. 7 (limitations).